\documentclass[a4paper,10pt,twocolumn]{article}

\usepackage[UTF8,fontset=fandol]{ctex}
\usepackage{amsmath,amssymb}
\usepackage{graphicx}
\usepackage{booktabs}
\usepackage{geometry}
\usepackage{hyperref}
\usepackage{caption}
\usepackage{array}
\usepackage{float}
\usepackage{xcolor}

\geometry{a4paper,margin=1.7cm,columnsep=0.75cm}
\setlength{\parindent}{2em}
\setlength{\parskip}{0.25em}
\renewcommand{\arraystretch}{1.15}

\title{课程报告三：基于条件 GAN 的小样本服饰图像生成}
\author{方俊睿\\学号：3230104069}
\date{}

\begin{document}
\maketitle

\noindent 仓库地址：\url{https://github.com/Tomabsolute/Neul.git}

\section{实验背景介绍}
生成对抗网络（Generative Adversarial Network, GAN）通过生成器与判别器的对抗训练学习数据分布。生成器从随机噪声中合成样本，判别器判断输入是真实样本还是生成样本；两者在交替优化中形成博弈，使生成器逐步产生更接近真实数据的图像。

普通 GAN 只能进行无条件生成，无法直接控制输出类别。为了增强可控性，本实验选择条件 GAN（Conditional GAN, CGAN）作为作业选题，将类别标签作为额外条件同时输入生成器和判别器，从而实现“指定类别生成图像”。与大规模人脸或复杂自然图像生成相比，Fashion-MNIST 尺寸小、类别清晰、训练成本低，更适合作为课程作业中的小样本 GAN 实验。

本实验代码参考了 \texttt{work3/demo-gan/} 中的 MNIST GAN 与 CGAN notebook，并整理为可复现的脚本工程，包含训练脚本、推理脚本、结果可视化与模型下载说明。

\section{实验任务与目标}
本实验目标是在小样本 Fashion-MNIST 上训练一个轻量级条件 DCGAN，实现以下功能：
\begin{itemize}
  \item 使用类别标签控制生成结果，如指定生成 Trouser、Sneaker、Bag 等类别；
  \item 提供完整训练脚本 \texttt{train\_cgan.py}，可自动下载数据集并保存 checkpoint；
  \item 提供推理脚本 \texttt{infer.py}，可加载训练好的生成器并输出图像网格；
  \item 控制数据量与训练轮次，不追求高精度，只要求能观察到 GAN 从噪声到类别轮廓的学习过程。
\end{itemize}

\section{数据集与预处理}
\subsection{数据集介绍}
实验使用 Fashion-MNIST 数据集。该数据集包含 10 类灰度服饰图像，每张图像大小为 $28\times28$，类别包括 T-shirt/top、Trouser、Pullover、Dress、Coat、Sandal、Shirt、Sneaker、Bag 与 Ankle boot。

为满足小样本训练要求，默认只从训练集中随机采样 6000 张图像。该设置既能保证每个类别有一定样本数量，又能将训练时间控制在较短范围内。

\subsection{预处理}
图像预处理较简单：
\begin{itemize}
  \item 转为张量；
  \item 将像素从 $[0,1]$ 归一化到 $[-1,1]$；
  \item 生成器输出层使用 \texttt{Tanh}，与该归一化范围保持一致。
\end{itemize}

\section{模型结构}
\subsection{生成器}
生成器输入由两部分组成：随机噪声 $z$ 与类别标签嵌入 $e(y)$。二者拼接后先经过全连接层映射到 $128\times7\times7$ 的特征图，再通过两层反卷积逐步上采样到 $1\times28\times28$。

\begin{table}[H]
\centering
\small
\caption{生成器结构}
\begin{tabular}{c|c}
\toprule
模块 & 输出尺寸 \\
\midrule
Noise + Label Embedding & $100+50$ \\
Linear + BN + ReLU & $128\times7\times7$ \\
ConvTranspose + BN + ReLU & $64\times14\times14$ \\
ConvTranspose + Tanh & $1\times28\times28$ \\
\bottomrule
\end{tabular}
\end{table}

\subsection{判别器}
判别器同样接收类别标签。具体做法是将标签嵌入映射为 $28\times28$ 的条件图，再与真实或生成图像在通道维拼接，形成 $2\times28\times28$ 输入。随后通过两层卷积与全连接层输出真假 logit。

\begin{table}[H]
\centering
\small
\caption{判别器结构}
\begin{tabular}{c|c}
\toprule
模块 & 输出尺寸 \\
\midrule
Image + Label Map & $2\times28\times28$ \\
Conv + LeakyReLU & $64\times14\times14$ \\
Conv + BN + LeakyReLU & $128\times7\times7$ \\
Flatten + Linear & $1$ \\
\bottomrule
\end{tabular}
\end{table}

\section{训练策略}
实验采用经典 GAN 交替训练方式。每个 batch 先更新判别器，再更新生成器。损失函数为 \texttt{BCEWithLogitsLoss}，因此判别器最后一层不额外接 \texttt{Sigmoid}。

\begin{align}
\mathcal{L}_D &= -\mathbb{E}_{x,y}[\log D(x,y)]-\mathbb{E}_{z,y}[\log(1-D(G(z,y),y))],\\
\mathcal{L}_G &= -\mathbb{E}_{z,y}[\log D(G(z,y),y)].
\end{align}

\begin{table}[H]
\centering
\small
\caption{默认训练配置}
\begin{tabular}{c|c}
\toprule
项目 & 配置 \\
\midrule
数据集 & Fashion-MNIST \\
训练样本数 & 6000 \\
Epoch & 20 \\
Batch size & 128 \\
噪声维度 & 100 \\
标签嵌入维度 & 50 \\
优化器 & Adam \\
学习率 & $2\times10^{-4}$ \\
Adam $\beta_1$ & 0.5 \\
设备 & CUDA / MPS / CPU 自动选择 \\
\bottomrule
\end{tabular}
\end{table}

\section{代码与运行方式}
项目主要文件如下：
\begin{itemize}
  \item \texttt{work3/train\_cgan.py}：训练条件 DCGAN；
  \item \texttt{work3/infer.py}：加载 checkpoint 并生成图片；
  \item \texttt{work3/src/models.py}：生成器、判别器与 DCGAN 初始化；
  \item \texttt{work3/visualize\_results.py}：根据 \texttt{history.csv} 绘制曲线；
  \item \texttt{work3/download\_model.py}：从开源平台或直链下载模型文件。
\end{itemize}

一键运行命令为：
\begin{verbatim}
bash work3/run_all_experiments.sh
\end{verbatim}

单独训练命令为：
\begin{verbatim}
python3 work3/train_cgan.py \
  --dataset fashion-mnist \
  --download \
  --epochs 20 \
  --batch-size 128 \
  --num-samples 6000
\end{verbatim}

推理命令为：
\begin{verbatim}
python3 work3/infer.py \
  --checkpoint work3/results/cgan_fashion_mnist/checkpoints/best_generator.pt \
  --labels 0,1,2,3,4,5,6,7,8,9 \
  --repeat 8 \
  --output work3/results/figures/inference_grid.png
\end{verbatim}

\section{模型保存与下载}
训练脚本会保存：
\begin{itemize}
  \item \texttt{best\_generator.pt}：用于最终推理的生成器 checkpoint；
  \item \texttt{last\_checkpoint.pt}：最后一轮完整 checkpoint；
  \item \texttt{history.csv}：每个 epoch 的损失与判别器分数；
  \item \texttt{samples/epoch\_*.png}：训练过程中的固定噪声生成结果。
\end{itemize}

提交到 Gitee 或 GitHub 时，可将 \texttt{best\_generator.pt} 作为 Release 附件上传。下载后使用如下命令恢复模型：
\begin{verbatim}
python3 work3/download_model.py \
  --url "模型文件直链" \
  --output work3/results/cgan_fashion_mnist/checkpoints/best_generator.pt
\end{verbatim}

\section{实验结果记录方式}
GAN 没有像分类任务一样直接的准确率指标，因此本实验主要记录以下训练信号：
\begin{itemize}
  \item 判别器损失 \texttt{d\_loss}；
  \item 生成器损失 \texttt{g\_loss}；
  \item 判别器对真实样本的平均分数 \texttt{D(real)}；
  \item 判别器对生成样本的平均分数 \texttt{D(fake)}；
  \item 固定噪声与固定标签下的生成图像变化。
\end{itemize}

理想情况下，训练初期生成图像接近随机噪声；训练若干轮后，不同标签对应的整体轮廓会逐渐分离，例如 Trouser 呈现竖直长条结构，Bag 更接近块状轮廓，Sneaker 与 Ankle boot 呈现鞋类横向轮廓。由于默认训练量较小，局部纹理不一定清晰，但应能观察到类别条件约束带来的可控生成效果。

\section{结果分析}
\subsection{条件信息的作用}
与无条件 GAN 相比，CGAN 在生成器和判别器两端都加入标签条件，使模型学习 $p(x|y)$ 而非单一的 $p(x)$。生成器因此可以根据标签改变输出分布；判别器也不只是判断图像真假，还要判断图像与标签是否匹配。这种设置让推理阶段可以通过标签控制生成类别。

\subsection{小样本训练的影响}
默认只使用 6000 张图像，相当于完整训练集的十分之一。小样本设置降低了训练成本，但也更容易带来模式坍塌和类别混淆。例如 Shirt 与 T-shirt/top、Coat 与 Pullover 本身形状接近，在训练轮次较少时生成结果可能边界不明显。

\subsection{损失曲线解释}
GAN 的损失不应像分类任务一样单调下降。若判别器过强，\texttt{D(real)} 接近 1 且 \texttt{D(fake)} 接近 0，生成器梯度可能变弱；若判别器过弱，生成器可能获得虚假的低损失。较合理的状态是二者保持动态平衡，固定噪声生成图像逐步稳定并呈现类别轮廓。

\section{实验不足与改进方向}
\begin{itemize}
  \item 只使用低分辨率灰度图像，生成结果主要体现轮廓，视觉复杂度有限；
  \item 未引入 FID、IS 等生成质量指标，评价主要依赖曲线与可视化；
  \item 小样本训练下可能出现模式坍塌，可进一步加入谱归一化、标签平滑或 WGAN-GP；
  \item 如果希望生成更高质量图像，可迁移到 CIFAR-10 或 CelebA，并使用更深的 DCGAN / StyleGAN 结构。
\end{itemize}

\section{结论}
本实验完成了一个可复现的基于 GAN 的生成模型任务：使用条件 DCGAN 在小样本 Fashion-MNIST 上进行类别可控图像生成。项目提供了训练脚本、推理脚本、可视化脚本与模型下载脚本，满足“代码可获得、包含 train 脚本、模型可下载并运行推理”的课程提交要求。虽然默认训练配置不追求高精度和高画质，但能够展示 GAN 对抗训练与条件生成的核心思想。

\end{document}
